\section{Projective examples with different nilpotent structures}
\label{sec:examples}

We construct two partial resolutions of the same anticanonical
hypersurface. Their deformation spaces have smooth reductions of
dimension six and Zariski tangent spaces of dimension eight. The
nilradical has square zero in the first example and nonzero square in
the second.

Let $N=\{y\in\Z^5:\sum_i y_i=0\}$ and take the reflexive polytope
\[
  \Pi=[-1,1]^5\cap N_{\mathbb R}.
\]
Its thirty vertices are the permutations of $(1,1,-1,-1,0)$.
Fixing two coordinates with equal signs gives a unimodular triangle.
Fixing $y_i=1$, $y_j=-1$ gives a hexagon $H_{ij}$ whose remaining
coordinates are the permutations of $(1,0,-1)$. These are all the
two-dimensional faces. Each of the twenty hexagons has one interior
lattice point, its centre, and primitive edges. The polar has vertices
$\pm\bar e_i$ in $\Z^5/\Z(1,1,1,1,1)$, so the edge dual to each
hexagon is primitive.

Take a general anticanonical hypersurface in the toric fourfold
defined by the face fan. For a set $T$ of hexagons let $X_T$ be its
proper transform after star subdivision at the centres of all hexagons
outside $T$. These projective crepant subdivisions leave one exact
degree-six cone at each member of $T$. Indeed, the hypersurface meets
its one-dimensional orbit transversely in a single point, and the
transverse toric threefold is the cone of the marked hexagon. The unit
triangles and primitive edges give smooth germs elsewhere. Inserting
the remaining centres and taking a coherent maximal crepant
triangulation refining the resulting fan gives the required smooth
projective resolution $Y$. Its local restriction is the vertex blow-up
at each remaining cone \cite[\S8.1]{W4}.
The primal and polar lattice-point counts are $51$ and $11$,
respectively. Neither has a lattice point interior to a facet.
The dual-edge corrections in Batyrev's Hodge number formula vanish,
so $(h^{1,1}(Y),h^{2,1}(Y))=(46,6)$
\cite[Corollary~4.5.1]{Batyrev}.
The direct-image and higher-direct-image calculations for the tangent
sheaf at the vertex blow-up give $h=h^{2,1}(Y)=6$,
as in \cite[proof of Lemma~5.2]{W5}.

We take the ordered triples
\[
 A=(H_{54},H_{53},H_{51}),\qquad
 B=(H_{54},H_{24},H_{14}),\qquad
 C=(H_{24},H_{23},H_{21}).
\]
Thus $A$ and $B$ have the common face $H_{54}$, whereas $A$ and $C$
are disjoint. At each $H_{ij}$ read the free coordinates in increasing
order of their indices. Order its vertices $v_0,\ldots,v_5$ by the
successive triples
\[
 (-1,0,1),\ (0,-1,1),\ (1,-1,0),\ (1,0,-1),\ (0,1,-1),\ (-1,1,0).
\]
Write $D_k$ for the intersection vector of the $k$th boundary curve
against the ambient prime divisors. If $e_v$ denotes the coordinate of
the ray through $v$ and $c$ is the centre, the wall relation gives
\[
 D_k=e_{v_{k-1}}+e_{v_{k+1}}-e_{v_k}-e_c
 \qquad(k\bmod 6).
\]
Use the three rows $D_0-D_3,D_4-D_0,D_2-D_4$ for the coordinates
$(s,u,v)$. The first spans the line and the last two span the plane.
By \cite[Proposition~8.3]{W4} the ambient degrees detect all homology
relations here. This also follows from the vanishing correction term in
$h^{1,1}$, since the invariant divisors then span $H^2(Y;\Q)$.

\begin{proposition}\label{prop:cube-examples}
Put $R=\C\{w_1,\ldots,w_6\}$. The analytic deformation rings of
$X_{A\cup B}$ and $X_{A\cup C}$ are respectively
\[
 R\{x,y\}/(x^2,xy,y^2)
 \qquad\text{and}\qquad
 R\{x,y\}/(x^2,y^2).
\]
Every deformation over a reduced base is locally trivial at the
marked vertices in both examples.
\end{proposition}

\begin{proof}
The following table gives the unique relation, normalized by its first
entry, on each ordered triple. Entries are the $(s,u,v)$ coefficients
on successive faces. Extend each relation by zero on the other faces.
\[
\begin{array}{c|rrr|rrr|rrr}
 &\multicolumn{3}{c|}{1}&\multicolumn{3}{c|}{2}&\multicolumn{3}{c}{3}\\
 A &1&0&-1&-1&0&1&-1&-1&-1\\
 B &1&1&0&1&0&1&-1&0&-1\\
 C &1&-1&0&-1&1&0&1&-1&0
\end{array}
\]
The left kernels for $A\cup B$ and $A\cup C$ have the indicated
pairs of rows as bases. More explicitly, after dropping the centre
columns, which cancel, the five-face matrix has $15$ rows of rank
$13$, and the six-face matrix has $18$ rows of rank $16$. All entries
are given by the wall formula above.

On $A\cup B$ let $x,y$ be the coefficients of the two displayed
relations. The face $H_{51}$ gives $x^2=0$, and $H_{24}$ gives
$y^2=0$. The common face $H_{54}$ has block $(x+y,y,-x)$, giving
$xy=0$. On $A\cup C$, the two relations have disjoint support and
each has a nonzero product of a line and a plane coefficient on one
of its faces. The restricted quadrics thus span precisely $x^2,y^2$.
Theorem~\ref{thm:main} gives the two rings. In each case their
reductions are defined by $x=y=0$, where all local maps are constant.
\end{proof}

The nilradical in the second ring has cube zero and its square is
the free $R$-module generated by $xy$. Its successive quotients,
including the reduction, have ranks $(1,2,1)$ over $R$. The
corresponding ranks for the first ring are $(1,2)$. In particular, the distinction
between the two unrestricted germs persists after any analytic change
of coordinates.

\subsection{Two cones in the presence of original nodes}
For comparison consider the following partial resolutions, also obtained
by inserting selected hexagon centres. For the first, take in
$\Z\oplus\Z^3$ the convex hull of $(1,0)$, the six points
$(0,\pm e_i)$, and the sixteen points $(-1,a)$ where $a$ runs through
the permutations of $(\pm1,\pm1,0)$ and the sign triples with
$a_1a_2a_3=1$. Retain the hexagons with centres $(-1,e_1)$ and
$(-1,e_2)$. For the second take $X_+$ of \cite[\S5]{Partial},
with the two markings specified there. The two polytopes are entries
248 and 528 in the census accompanying \cite{W5}.

In both cases all primal edges and all dual edges of singular faces
are primitive. Every triangle is unimodular, every quadrilateral is
a unit square, and the remaining two-faces are degree-six hexagons.
The same coherent refinement as above resolves the retained vertices
and selects a small resolution at every square. The following counts
and ranks use the two retained hexagons in each case. Here $N_0$ is
the matrix of square circuits $e_i+e_j-e_k-e_l$, for opposite pairs
$(i,j),(k,l)$, and $M$ adjoins
the three rows for $(s,u,v)$ from each retained hexagon.
\[
\begin{array}{c|rrr|rr|rr}
 &\text{triangles}&\text{squares}&\text{hexagons}&
 h^{1,1}(Y)&h^{2,1}(Y)&\rk N_0&\rk M\\ \hline
 248&52&12&6&25&25&12&17\\
 528&32&30&4&27&18&20&22
\end{array}
\]
The accompanying computation reconstructs the faces from the ordered
vertices, checks unimodularity and the Hodge counts, and forms the
curve matrices from the wall formula and square circuits. The input
files contain only vertices and cyclic hexagon orders.

For the first threefold the unique relation has cone coefficients
$(1,1,1;-1,-1,-1)$. No choice of local branches contains this
direction. The reduced base is smooth of dimension $25$, and the
unrestricted ring is $\C\{w_1,\ldots,w_{25},t\}/(t^2)$ by
Theorem~\ref{thm:main}. In particular the nilpotent direction is
free over the entire reduced base.

For $X_+$ the kernel on the original nodes has dimension ten.
In the cyclic markings of \cite[\S5]{Partial}, the image on the two
cones has precisely the equations $u_0=0$ and $v_0+s_1=0$.
Writing $t=s_1$, $x=s_0$, $y=u_1$, $z=v_1$, the unrestricted ring is
\[
 \C\{w_1,\ldots,w_{28},t,x,y,z\}/(tx,ty,tz).
\]
The $28$ free coordinates include the ten original-node relations.
The two components have dimensions $31$ and $29$ and intersect along
the smooth germ of dimension $28$. A general point of the
$31$-dimensional component gives a smoothing. A general fibre on the other
component has one original node and one node near the first cone. In the ordering of \cite[\S5]{Partial}, the marked incidence relations
force precisely the period of original node $17$ and one of the three
plane periods to vanish on this component.
